% Generated draft: requires mathematical and source review.
\documentclass[11pt]{article}
\usepackage[margin=25mm]{geometry}
\usepackage{fontspec}
\setmainfont{Latin Modern Roman}
\usepackage{amsmath,amssymb,mathtools,unicode-math}
\setmathfont{Latin Modern Math}
\usepackage{enumitem,needspace,xurl,hyperref}
\setlength{\parindent}{0pt}
\setlength{\parskip}{5pt}
\setlength{\emergencystretch}{3em}
\newcommand{\N}{\mathbb{N}}
\newcommand{\Z}{\mathbb{Z}}
\newcommand{\Q}{\mathbb{Q}}
\newcommand{\R}{\mathbb{R}}
\newcommand{\C}{\mathbb{C}}
\newcommand{\F}{\mathbb{F}}
\newcommand{\A}{\mathbb{A}}
\newcommand{\PP}{\mathbb{P}}
\newcommand{\E}{\mathbb{E}}
\newcommand{\eps}{\varepsilon}
\newcommand{\dd}{\,\mathrm{d}}
\begin{document}
\section*{UnsolvedMath ID: 800011}
\textbf{Source problem code:} catalogue entry
\subsection{Definitions and mathematical statement}
An ordered semigroup, under the usual partial-order convention, is a nonempty semigroup \((S,\cdot)\) equipped with a partial order \(\leq\) such that, for all \(a,b,c\in S\),
\[
a\leq b\ \Longrightarrow\ ca\leq cb\ \text{ and }\ ac\leq bc.
\]
Under the inequational convention, a variety of ordered semigroups is the class of all ordered semigroups satisfying a specified collection of universally quantified inequalities \(u\leq v\), where \(u\) and \(v\) are nonempty words in variables. Let \(\mathcal L\) be the lattice of these varieties, ordered by inclusion. A variety \(\mathcal W\in\mathcal L\) covers \(\mathcal V\in\mathcal L\) if
\[
\mathcal V\subsetneq\mathcal W
\quad\text{and}\quad
\nexists\mathcal U\in\mathcal L\;\bigl(\mathcal V\subsetneq\mathcal U\subsetneq\mathcal W\bigr).
\]
\textbf{Problem (E. S. Lyapin).} Describe all varieties of ordered semigroups that have a cover; equivalently, characterize
\[
\left\{\mathcal V\in\mathcal L:\ \exists\mathcal W\in\mathcal L\ \text{such that }\mathcal W\text{ covers }\mathcal V\right\}.
\]
The supplied statement reports that some such varieties are described in references [Mak2] and [Mak3], and that not every variety in this lattice has a cover, citing [Mak2]. These are supplied source assertions, not independently verified conclusions. The record does not explicitly define ordered semigroups or their varieties; the partial-order and inequational conventions above require confirmation against the source.
\subsection{Short English statement}
Characterize precisely which varieties of ordered semigroups admit an immediate strictly larger variety in the inclusion lattice. The problem is attributed to E. S. Lyapin. The registry labels it open; its intended conventions and the incomplete references [Mak2] and [Mak3] require human review.
\subsection{Sources}
\begin{itemize}
\item Local UnsolvedMath registry, permanent ID \(800011\), source problem code \texttt{catalogue entry}, title \emph{The Lyapin's Notebook --- Problem 11}. The supplied record labels the problem open and records Bershadsky, Kublanovsky, and Mashevitzky in its proposer field, with year 2026; the mathematical statement itself attributes the problem to E. S. Lyapin.
\item Bershadsky, Kublanovsky, and Mashevitzky, \emph{The Lyapin's Notebook: Unsolved Problems in Semigroup Theory} (2026), Problem 11. Supplied source URL: \url{https://arxiv.org/abs/2604.04763}. The record reports access on 2026-10-16 and extraction by \texttt{arxiv\_latex}; this provenance is supplied and has not been independently verified here.
\item [Mak2], cited in the supplied statement for examples of varieties having covers and for the assertion that some varieties have no cover. No bibliographic details for this reference are supplied.
\item [Mak3], cited in the supplied statement for examples of varieties having covers. No bibliographic details for this reference are supplied.
\end{itemize}
\end{document}
